\section{A quantitative singular-set bound for weighted minimizers}
\label{sec:rpi:weighted-minimizers}

The smooth positive mass theorem used in
Section~\ref{sec:rpi:smooth-repair} is
\cite[Corollary~1.6]{BrendleWang2026}. In its dimension-descent
proof, a weighted perimeter minimizer with a volume forcing is
completed across its singular set. This appendix verifies the
quantitative covering estimate needed for that variational class.
It supplies that input to the cited proof; the positive mass
conclusion remains the external theorem.

The quantitative input here is
\cite[Theorem~1.3]{NaberValtorta2020}, with the Riemannian
bounded-mean-curvature hypotheses (1.14)--(1.15) of preprint
version~4. Bounded mean curvature alone does not place a singular
set in a quantitative stratum of codimension eight. We use local
minimization to prove that inclusion at every sufficiently small
scale, then apply the quantitative theorem.

\begin{lemma}[A quantitative estimate for weighted minimizers]
\label{lem:rpi:bw-quantitative}
Let $n\ge8$, let $g$ be smooth on an open set $U$, and let
$\widehat\rho>0$ and $\Phi$ be smooth there. Put
\[
 g_* =\widehat\rho^{2/(n-1)}g,
 \qquad f=\Phi\widehat\rho^{-1/(n-1)}.
\]
Suppose that a locally finite-perimeter set $E$ locally minimizes
\[
 \mathcal F(E;W)
   =P_{g_*}(E;W)-\int_{E\cap W}f\,dV_{g_*}
\]
against changes compactly supported in $W\Subset U$. Let
$\mathcal S$ be a compact subset of the singular set of its
perimeter support in $U$. There are constants $C,t_0>0$ such that
\begin{equation}\label{eq:rpi:bw-tube}
 \Vol_g\bigl(B_t^g(\mathcal S)\bigr)\le C t^8,
 \qquad 0<t<t_0.
\end{equation}
The constants are local: they may depend on the minimizer and on
a buffered compact neighborhood of $\mathcal S$.
\end{lemma}
\begin{proof}
The proof has three parts. Absorbing the density into the metric
gives bounded generalized mean curvature and strict blow-up
compactness. Compactness and low-dimensional regularity then exclude
cones with too many translation symmetries at every small scale.
The quantitative stratification theorem converts this exclusion into
the tube estimate.

The conformal metric absorbs the positive perimeter density:
$d\mathcal H_{g_*}^{n-1}=\widehat\rho\,
d\mathcal H_g^{n-1}$. Also
$f\,dV_{g_*}=\Phi\widehat\rho\,dV_g$, so the displayed
functional is exactly the corresponding weighted perimeter minus
weighted volume functional. All estimates below are first made
in $g_*$. Choose neighborhoods
\[
 \mathcal S\subset U_0\Subset U_1\Subset U.
\]
Smoothness and positivity make their metric and forcing bounds
uniform after working on a slightly larger compact neighborhood
of $\overline U_1$.

\emph{First variation and local estimates.}
Let $V=|\partial^*E|$ be the multiplicity-one boundary varifold,
computed in $g_*$. Domain variation by any compactly supported
smooth vector field $Y$ gives
\[
 \delta_{g_*}V(Y)
    =\int_{\partial^*E} f\langle Y,\nu_E\rangle_{g_*}\,
                      d\mathcal H_{g_*}^{n-1}.
\]
Here the divergence theorem for finite-perimeter sets is applied
to $fY$. Thus $V$ has bounded generalized mean curvature, with
no additional first-variation measure on its singular set. In
particular
\[
 |\delta_{g_*}V(Y)|
       \le\|f\|_{L^\infty(U_1)}\int |Y|_{g_*}\,d\|V\|.
\]
The minimizing comparison also gives
\begin{equation}\label{eq:rpi:bw-almost-minimizing}
 P_{g_*}(E;W)\le P_{g_*}(F;W)
       +\|f\|_{L^\infty(U_1)}
                       \Vol_{g_*}(E\triangle F)
\end{equation}
for changes supported in $W\Subset U_1$. The usual ball
comparisons and relative isoperimetric inequality for this
bounded-volume-forcing functional give uniform local perimeter
upper bounds and two-phase density lower bounds. In particular,
for $x$ in the perimeter support in $\overline U_0$ and all
sufficiently small $s$, both phases occupy at least $c s^n$ of
$B_s^{g_*}(x)$, and
\[
 c s^{n-1}\le P_{g_*}(E;B_s^{g_*}(x))\le C s^{n-1}.
\]
These are the local estimates for almost-minimizing boundaries;
they do not say that the boundary is Euclidean area minimizing
in a fixed coordinate chart.

\emph{Strict blow-up compactness.}
Consider singular points $x_j\in\overline U_0$ and radii
$r_j\downarrow0$. Choose $g_*$-orthonormal frames at $x_j$ and
write
\[
 \Psi_j(y)=\exp_{x_j}^{g_*}(r_jy),\qquad
 g_j=r_j^{-2}\Psi_j^*g_*,\qquad
 f_j(y)=r_j f(\Psi_j(y)).
\]
The rescaled sets $E_j=\Psi_j^{-1}(E)$ minimize
$P_{g_j}-\int f_j\,dV_{g_j}$ on domains exhausting $\R^n$.
On every fixed ball $g_j\to\delta$ and $f_j\to0$ smoothly.
The preceding perimeter bounds give, after passage to a
subsequence,
\[
 \chi_{E_j}\longrightarrow\chi_{E_\infty}
          \quad\hbox{in }L^1_{\rm loc}(\R^n).
\]
We need strict perimeter convergence as well, since $L^1$
convergence by itself could lose cancelling sheets.

Fix $0<a<b$ in a bounded coordinate ball. Coarea permits radii
$s_j\in(a,b)$ for which the traces exist, neither perimeter
measure charges $\partial B_{s_j}$, and
\[
 \int_{\partial B_{s_j}}
          |\chi_{E_j}-\chi_{E_\infty}|\,
                d\mathcal H^{n-1}\longrightarrow0.
\]
Glue $E_\infty$ inside $B_{s_j}$ to $E_j$ outside. The
modification is compactly supported in a larger fixed ball.
The BV gluing formula and minimality give
\begin{align*}
 P_{g_j}(E_j;B_{s_j})
 &\le P_{g_j}(E_\infty;B_{s_j})
       +C\int_{\partial B_{s_j}}
            |\chi_{E_j}-\chi_{E_\infty}|\,
                           d\mathcal H^{n-1}\\
 &\quad+\|f_j\|_{L^\infty(B_b)}
       \Vol_{g_j}\bigl((E_j\triangle E_\infty)
                                      \cap B_{s_j}\bigr).
\end{align*}
The last two terms tend to zero. Uniform convergence of the
metrics compares their perimeters with Euclidean perimeter by
factors $1+o(1)$, so
\[
 \limsup_jP_\delta(E_j;B_a)
             \le P_\delta(E_\infty;B_b).
\]
Let $b\downarrow a$ at a radius with
$|D\chi_{E_\infty}|(\partial B_a)=0$, and use lower
semicontinuity. It follows that
\begin{equation}\label{eq:rpi:bw-strict-bv}
 P_\delta(E_j;B_a)\longrightarrow
                         P_\delta(E_\infty;B_a).
\end{equation}
The same argument on other relatively compact balls gives local
perimeter-measure convergence. Together with
$D\chi_{E_j}\rightharpoonup D\chi_{E_\infty}$,
Reshetnyak continuity now gives boundary-varifold convergence:
the unoriented tangent plane is a continuous even function of
the measure-theoretic normal. Changing between the $g_j$
and Euclidean varifold conventions introduces only vanishing
metric errors. In particular there is no residual varifold mass
from disappearing or cancelling sheets.

The limit is a Euclidean locally perimeter-minimizing boundary.
Indeed, for $F\triangle E_\infty\Subset B_a$, repeat the
same gluing with $F$ inside $B_{s_j}$. On the gluing annulus
$F=E_\infty$, so the trace error is unchanged. Letting
$j\to\infty$ and then $b\downarrow a$ proves
$P(E_\infty;B_a)\le P(F;B_a)$ at continuity radii. Such a
ball can contain any compact modification. Finally, the two-phase
density bounds at $x_j$ pass to every fixed rescaled ball
centered at zero. Both phases of $E_\infty$ have positive
density there; in particular its boundary varifold is nonzero.

\emph{Uniform exclusion of highly symmetric cones.}
We claim that there are single constants $\epsilon,R_0>0$
such that, at every singular point in $\overline U_0$ and
every physical scale $0<s<R_0$, the rescaled varifold is not
$\epsilon$-close to an $(n-7)$-symmetric comparison cone.
Comparison is in the weak-varifold metric on $B_1$ used in
quantitative stratification, with balls contained in the buffered
neighborhood. Such a cone is dilation invariant about zero and
translation invariant along an $(n-7)$-plane; it is not assumed
to be a minimizing boundary.

If the claim fails, choose $x_j,r_j$ as above and comparison
cones $C_j$ with distance less than $1/j$ on $B_1$. Denote
their invariant $(n-7)$-planes by $K_j$ and pass to
$K_j\to K$ in the Grassmannian. Closeness on the open ball
gives a uniform mass bound for $C_j$ on $B_{1/2}$, by using
a cutoff equal to one there and supported in $B_{3/4}$.
No mass bound at the boundary of the comparison ball is needed.
Cone homogeneity propagates this interior bound to every fixed
$B_R$. Radon compactness on the Grassmann bundle therefore
gives a global varifold limit $C$, homogeneous about zero and
translation invariant along $K$. The strict compactness just
proved identifies it on the comparison ball:
\begin{equation}\label{eq:rpi:bw-cone-identification}
 C\llcorner B_1
    =|\partial^*E_\infty|\llcorner B_1.
\end{equation}
There was no first-variation bound for the arbitrary $C_j$.
Consequently it is this identification with the minimizing-boundary
limit, not their mass bounds alone, that gives rectifiability,
integrality and multiplicity one in $B_1$. Homogeneity extends
these properties to $C$ globally.

We next construct a boundary-set cone; we do not assert that the
moving-center limit $E_\infty$ is conical outside $B_1$.
The radial direction is tangent almost everywhere to a rectifiable
cone. Equation~\eqref{eq:rpi:bw-cone-identification} therefore
implies, on $B_1\setminus\{0\}$,
\[
 x\cdot D\chi_{E_\infty}=0,
 \qquad
 \divg(x\chi_{E_\infty})=n\chi_{E_\infty}\,dx.
\]
The second expression includes the dimensional term. In polar
coordinates the first identity says that the phase is constant
along almost every radial segment. Extend those phase values
radially to a conical set $E_0\subset\R^n$. It agrees with
$E_\infty$ almost everywhere in $B_1$. The homogeneous
perimeter bound gives locally finite perimeter also at the
origin: cutting out a ball of radius $a$ introduces at most
$O(a^{n-1})$ additional perimeter, which vanishes as
$a\downarrow0$. Hence
\[
 |\partial^*E_0|=C
 \quad\hbox{on }\R^n,
\]
first by agreement inside $B_1$, and then by homogeneity.

The conical set $E_0$ is globally locally perimeter minimizing.
Given a compact modification, homothetically contract a ball
containing its support into $B_1$. In that ball $E_0$ equals
$E_\infty$, so the minimizing comparison is available.
Undoing the contraction and using the $(n-1)$-degree scaling
of perimeter proves the comparison at the original scale. Only
the phase extension $E_0$, not the earlier limit outside
$B_1$, is used in this scaling argument.

\emph{Splitting as a minimizing boundary.}
Translation invariance of the rectifiable varifold $C$ along
$K$ makes every $v\in K$ tangent almost everywhere. Thus
\[
 v\cdot D\chi_{E_0}=0\qquad(v\in K).
\]
Distributional constancy in those directions and Fubini yield
$E_0=K\times F$ up to a null set, for a locally
finite-perimeter conical set $F\subset K^\perp$.
This preserves the phase information that would be absent from
a statement about splitting stationary varifolds alone.

The factor $F$ minimizes perimeter. To see this directly, let
$\ell=\dim K$ and $G\triangle F\Subset D\subset K^\perp$
be a bounded compact modification. Over a cube $Q_L\subset K$
of side $L$, replace $F$ by $G$, retaining $E_0$ elsewhere.
The product gluing formula and the minimality of $E_0$ give
\[
 0\le |Q_L|\bigl[P(G;D)-P(F;D)\bigr]
                   +P(Q_L)|G\triangle F|.
\]
There is no interface contribution at $\partial D$ because
the phases agree near it. The displayed lateral cost at
$\partial Q_L$ is essential. Divide by $|Q_L|$ and let
$L\to\infty$; its coefficient tends to zero. Thus
$P(F;D)\le P(G;D)$. Here $\ell=n-7\ge1$; if there
are more invariant directions the residual dimension only
decreases.

The residual ambient dimension is at most seven. The classical
regularity theorem for minimizing boundaries therefore makes
the minimizing cone $F$ regular at its vertex
\cite{SimonGMT}. A nontrivial
boundary cone smooth at its vertex is a hyperplane, so $F$
is a halfspace. Nontriviality follows from the two-phase density
bounds. Consequently $C$ is a multiplicity-one hyperplane
through zero.

The strict boundary-varifold convergence to this plane,
$g_j\to\delta$ smoothly, and $f_j\to0$ now give
flat-boundary regularity in a fixed smaller ball. One can use
the metric-uniform almost-minimizing-boundary theorem or
Riemannian Allard regularity; the prescribed-mean-curvature
equation then bootstraps the graph to smoothness. This contradicts
singularity of $E_j$ at zero. The argument uses Riemannian
regularity with converging smooth metrics, not an unsupported
bound for Euclidean mean curvature obtained by changing
coordinates. It proves the claimed uniform $\epsilon,R_0$,
at every singular center in the patch and every scale below
$R_0$.

\emph{The complete quantitative scale interval.}
Choose finitely many balls $B_{\rho_i/2}^{g_*}(p_i)$
covering $\mathcal S$, with
\[
 B_{4\rho_i}^{g_*}(p_i)\Subset U_0,
 \qquad 4\rho_i<R_0.
\]
They may be shrunk so that the rescaled ambient geometry has
the sectional-curvature and injectivity-radius bounds in the
cited Riemannian quantitative theorem. The first-variation
bound and local perimeter bound already give its bounded
generalized mean curvature and finite mass assumptions.
The constants are finite on this fixed finite cover; no
uniformity over unrelated metrics is asserted.

Apply \cite[Theorem~1.3]{NaberValtorta2020} to
$B_{2\rho_i}^{g_*}(p_i)$ after rescaling lengths by
$\rho_i^{-1}$. Its target ball $B_1$ is
$B_{\rho_i}^{g_*}(p_i)$. For every singular point in that
target ball and every rescaled scale $0<s<1$, the comparison
ball lies in $B_{2\rho_i}^{g_*}(p_i)$ and has physical
radius $\rho_i s<R_0$. Thus the same symmetry exclusion
holds on the entire interval $[r,1)$, for each $0<r<1$.
In the notation of quantitative stratification, the rescaled
singular set in the target ball lies in
\[
 S_{\epsilon,r}^{\,n-8}
       \quad\hbox{for every }0<r<1.
\]
The stratum index is $n-8$ because the excluded cones have
$n-7$ invariant directions. In particular this statement is
stronger than a statement merely about tangent cones or a
sequence of scales.

The quantitative theorem gives rescaled tube volume at most
$C_i r^8$. Returning to physical variables, for
$0<t<\rho_i/2$ we obtain
\[
 \Vol_{g_*}\left(
 B_t^{g_*}\bigl(\mathcal S\cap
                    B_{\rho_i/2}^{g_*}(p_i)\bigr)\right)
       \le C_i\rho_i^{n-8}t^8.
\]
The entire displayed tube lies in $B_{\rho_i}^{g_*}(p_i)$,
which is the target ball where the theorem controls its volume.
For $t<\min_i\rho_i/2$, these finitely many tubes cover
$B_t^{g_*}(\mathcal S)$. Summing gives its $Ct^8$ bound.
Finally, smooth metric equivalence on the buffered compact
neighborhood compares small tubes and volume elements for $g$
and $g_*$, proving Equation~\eqref{eq:rpi:bw-tube}.
Ambient dimension eight corresponds to stratum index zero and
is included.
\end{proof}

The estimate supplies the covering bound used in
\cite[Theorem~3.34]{BrendleWang2026} for the weighted minimizers of
its Definition~3.19 and Lemma~3.20. Indeed, disjoint equal-radius
balls of radius $t$ centered on a fixed compact singular set number
at most $Ct^{8-n}$: each has volume at least $ct^n$, and their
union lies in the tube just estimated. A maximal disjoint family
has doubled balls covering the singular set. The resulting bound
also implies every estimate with an arbitrarily small loss in its
exponent, as required in that dimension-descent construction.
